% 原书分章 PDF 第 20—21 页：第 25.5.1 节。
\subsection{在目标参数的不同路径可微逼近类之间选择}

令 $\{\Psi_{h,j}(P):h\}$ 为以调节参数 $h>0$ 逼近 $\Psi(P)$ 的一族目标参数，
由离散选择 $j=1,\ldots,J$ 索引。
例如，$\Psi_{h,j}(P)$ 可以是前几节分析过的 $\Psi(P)$ 的核平滑逼近，
其核选择为与不超过 $j$ 次的多项式正交。
前几节预先将 $j$ 选得较大，但在实践中，不同 $j$ 取值的实际表现可能相差很大。
对于任意 $h,j$，令 $\psi^*_{h,j,n}$ 为 $\Psi_{h,j}(P_0)$ 的 CV-TMLE。
令 $\sigma_{h,j,n}^2$ 为 $\Psi_{h,j}$ 在 $P_0$ 处的典范梯度
$D_{h,j}^*(P_0)$ 的方差估计量，使对序列 $h_n\to0$，有
\[
 (\sigma^{-1}_{h_n,j,n}n^{0.5})
 (\psi_{h_n,j,n}-\Psi_{h_n,j}(P_0))\Rightarrow_dN(0,1).
\]
令 $h_n(j)$ 为上述类型的数据自适应选择器，
并假定正则条件成立，使得对于 $j=1,\ldots,J$，
$h_n(j)$ 最优地权衡 $\psi^*_{h,j,n}$ 相对于 $\psi_0$ 的方差与偏差。
令 $\psi^*_{j,n}\equiv\psi^*_{h_n(j),j,n}$，
$\sigma_{j,n}=\sigma_{h_n(j),j,n}$。
如何在由 $j$ 索引的不同逼近策略之间选择？
由于每个 $h_n(j)$ 都选得使 $\psi^*_{j,n}$ 的偏差平方与其方差同阶，
所有这些估计量 $\psi^*_{j,n}$ 的 MSE 都与 $\sigma_{j,n}^2/n$ 同阶。
因此，一个合理的 $j$ 选择器 $j_n$，可定义为使估计量的估计方差最小的取值：
\[
 j_n\equiv\argmin_j\sigma_{j,n}^2.
\]
这一选择保证，我们以趋于 1 的概率，选出使
$\psi^*_{j,n}$ 以最快速率收敛到 $\Psi(P_0)$ 的参数族。
这可视为针对路径可微目标参数的经验效率最大化方法
（empirical efficiency maximization）的推广（Rubin 和 van der Laan，2008）。

为进行统计推断，对于每个 $j$，我们会选择 $h_n(j)$ 的略微欠平滑版本，
从而可以忽略渐近偏差。
例如，在数据自适应选择器 $h_n(j)$ 的定义中，
把常数 $C_n$ 选为以 $\log n$ 这样的缓慢速率趋于无穷。
用 $h_n(j)$ 表示这样的选择。
进行统计推断时，如果把 $j_n$ 当成预先指定的值，
直接选择基于 $\psi^*_{j_n,n}$ 的欠平滑版本的置信区间，
就可能产生过于乐观的推断，因为 $j_n$ 本身是为最小化置信区间宽度而选出的。
我们提出如下置信区间调整。有
\[
\begin{aligned}
 Z_n(j)&\equiv\sigma_{j,n}^{-1}n^{0.5}
 (\psi^*_{j,n}-\Psi_{h_n(j),j}(P_0))\\
 &=\frac{1}{\sqrt n}\sum_{i=1}^{n}
 \sigma_{j,n}^{-1}D^*_{h_n(j),j}(P_0)+o_P(1),
 \qquad j=1,\ldots,J.
\end{aligned}
\]
因此，在正则条件下，有
$Z_n=(Z_n(j):j)\Rightarrow N(0,\Sigma_0)$，其中
\[
 \Sigma_0(j_1,j_2)=\lim_{n\to\infty}
 P_0\sigma^{-1}_{j_1,n}\sigma^{-1}_{j_2,n}
 D^*_{h_n(j_1),j_1}(P_0)D^*_{h_n(j_2),j_2}(P_0)
\]
是估计量 $\psi^*_{j_1,n}$ 与 $\psi^*_{j_2,n}$ 的标准化版本之间的渐近相关。
这里，$\sigma_{0,h,j}^2\equiv P_0\{D^*_{h,j}(P_0)\}^2$
表示 $D^*_{h,j}(P_0)$ 的真实方差，并且我们隐含地假定，
该协方差极限几乎处处存在。注意，
$\Sigma_0(j,j)=1$，$j=1,\ldots,J$。
令 $q_{0,95}$ 为 $N(0,\Sigma_0)$ 的绝对值之最大值的 0.95 分位数。
于是，$\psi^*_{j,n}\pm q_{0.95}\sigma_{j,n}/n^{0.5}$，
$j=1,\ldots,J$，构成 $\psi_0$ 的同时渐近 0.95 置信区间。
因此，$\psi^*_{j_n,n}\pm q_{0.95}\sigma_{j,n}/n^{0.5}$
是 $\psi_0$ 的 0.95 置信区间。
将 $q_{0.95}$ 定义中的 $\Sigma_0$ 替换为相合估计量 $\Sigma_n$，
即可相合地估计 $q_{0.95}$。
